\begin{textsample}{NEG Dim 4 – 2011\_c4 – Score -1.00 – 2011\_c4/enade\_2011\_ci\_comp\_dab220f8e3.txt}  \label{ex:f4_neg_121}
CAPÍTULO 4 PERCEPÇÃO DA PROVA As análises feitas neste capítulo tratam das percepções dos concluintes da Área de Computação sobre a prova aplicada no ENADE/2011.
Estas percepções foram mensuradas por meio de nove questões que avaliaram desde o grau de dificuldade da prova até o tempo gasto para concluí -la.
As percepções sobre a prova foram relacionadas com o desempenho dos estudantes e com a Grande Região de funcionamento do curso.
O questionário de percepção da prova encontra -se ao final do Anexo IV, que traz a reprodução da prova.
O desempenho dos estudantes foi classificado em quatro quartos.
Para tanto, esse desempenho foi ordenado de forma ascendente.
O percentil 25, P25, também conhecido como primeiro quartil, é a nota de desempenho que deixa um quarto ( 25 \% ) dos valores observados abaixo e três quartos acima.
A Figura 1 apresenta uma ilustração deste conceito.
O quarto inferior de desempenho é composto pelas notas abaixo do primeiro quartil.
Já o percentil 75, P75, também conhecido como terceiro quartil, é o valor para o qual há três quartos ( 75 \% ) dos dados abaixo e um quarto acima dele.
O quarto superior de desempenho é composto pelas notas iguais ou acima do terceiro quartil.
O percentil 50, P50, também conhecido como mediana, é o valor que divide as notas em dois conjuntos de igual tamanho.
O segundo quarto inclui valores entre o primeiro quartil ( P25 ) e a mediana.
O terceiro quarto contém os valores entre a mediana ( P50 ) e o terceiro quartil ( P75 ).
Vale ressaltar que percentis, quartis e medianas são pontos que não obrigatoriamente pertencem ao conjunto original de dados, ao passo que os quartos são subconjuntos dos dados originais.
A seguir, serão apresentados gráficos com resultados selecionados, relativos às nove questões avaliadas por grupos de estudantes.
Os gráficos apresentam nas barras o percentual de alunos que assinalaram uma das opções ou a soma das porcentagens daqueles que assinalaram duas ( ou três ) delas.
Por exemplo, para as questões 1 e 2, os gráficos apresentam a porcentagem total de participantes que assinalaram as opções ( D ) difícil e ( E ) muito difícil.
Em cada barra foram assinalados também os extremos do intervalo de confiança de 95 \% como linhas verticais unidas por uma linha horizontal na forma da letra H maiúscula, semelhantemente aos gráficos do Capítulo 3.
As Tabelas no Anexo II apresentam os valores absolutos e a distribuição percentual das alternativas válidas das nove questões, segundo o mesmo recorte de desempenho dos alunos e Grande Região de funcionamento do curso. 4.1 GRAU DE DIFICULDADE DA PROVA 4.1.1 Componente de Formação Geral Ao avaliarem “ Qual o grau de dificuldade desta prova na parte de Formação Geral? ” ( Questão 1 ), 27,3 \% do grupo de inscritos e presentes optaram pelas alternativas difícil ou muito difícil.
Entretanto, para mais da metade dos estudantes ( 54,9 \% ), o Componente de Formação Geral da prova foi considerado com grau de dificuldade médio ( Gráfico 4.1, Gráfico 4.2 e, no Anexo II, a Tabela II.1 ).
O percentual de estudantes que consideraram a prova como difícil ou muito difícil foi maior na região Centro-Oeste, onde a proporção foi de 29,7 \%, enquanto a de menor incidência foi a Nordeste, com 22,6 \%.
No Gráfico 4.1 é possível observar que esta diferença é estatisticamente significativa.
Nas Grandes Regiões, a proporção de presentes à prova que consideraram o Componente de Formação Geral como sendo de grau de dificuldade médio esteve entre 51,8 \% na região Centro-Oeste e 58,6 \% na região Norte.
O percentual de alunos que consideraram a prova difícil ou muito difícil \textbf{decresce} em função dos quartos de desempenho e a diferença entre eles é estatisticamente significativa.
No quarto de maior desempenho a proporção de alunos que consideraram a prova difícil ou muito difícil foi de 18,1 \% e no de desempenho inferior esta proporção foi de 35,5 \%.
Para todos os quartos de desempenho a alternativa modal para esta pergunta foi médio.
Nos quartos de desempenhos extremos 52,1 \% e 55,6 \% dos respondentes, respectivamente no primeiro e último quartos, afirmaram que o grau de dificuldade foi médio. 4.1.2 Componente de Conhecimento Específico Ao responderem à Questão 2 – “ Qual o grau de dificuldade desta prova na parte de Componente Específico? ” – 54,6 \% do grupo de estudantes classificaram -na como difícil ou muito difícil.
Destaca -se que esta proporção foi bem mais elevada do que em relação ao Componente de Formação geral ( 27,3 \% na Questão 1 ).
Além disso, o Componente de Conhecimento Específico da prova foi considerado com grau de dificuldade médio por menos da metade dos alunos, 40,3 \%. ( Gráfico 4.3, Gráfico 4.4, e, no Anexo II, a Tabela II.2 ).
A análise das respostas dos estudantes quanto ao grau de dificuldade do Componente de Conhecimento Específico da prova, agregado por Grande Região, mostra que a diferença entre a maior e a menor proporção de alunos que a avaliaram como difícil ou muito difícil é estatisticamente significativa : 57,9 \% na região Sul e 47,8 \% na Nordeste.
Apenas na região Nordeste a classe modal foi a opção médio, em todas as demais Grandes Regiões a escolha da opção difícil foi a de maior frequência para o Componente de Conhecimento Específico.
Considerando -se a avaliação da dificuldade das questões do Componente de Conhecimento Específico da prova, de acordo com o desempenho dos estudantes, se observa diferença estatisticamente significativa entre o percentual do quarto inferior ( 49,2 \% ) e os demais.
Nos outros três quartos, a proporção dos que classificaram a parte específica como difícil ou muito difícil variou de 55,2 \% ( último quarto ) a 58,4 \% ( terceiro quarto ).
A alternativa modal para a Questão 2 foi o grau médio para o quarto inferior com 43,3 \% das escolhas.
Para os outros quartos a opção difícil foi a mais escolhida com 45,0 \% dos estudantes do segundo quarto e 45,5 \% no quarto superior. 4.2 EXTENSÃO DA PROVA EM RELAÇÃO AO TEMPO TOTAL Indagados quanto à extensão da prova, em relação ao tempo total oferecido para a sua resolução ( Questão 3 ), os estudantes apontaram, com maior incidência, a alternativa que considerava a extensão adequada, para todas as agregações consideradas ( Gráfico 4.5, Gráfico 4.6, e, no Anexo II, a Tabela II.3 ).
O percentual de alunos que responderam ser a extensão da prova adequada foi de 62,5 \%.
Já 32,2 \% dos inscritos presentes consideraram que o exame foi longo ou muito longo e 5,3 \% o avaliaram como curto ou muito curto.
Entre as Grandes Regiões a proporção daqueles que avaliaram a prova como longa ou muito longa em relação ao tempo total destinado à sua resolução variou de 27,3 \% na região Norte a 34,4 \% na região Nordeste.
A diferença entre o menor percentual, obtido na região Norte, é estatisticamente significativa em relação aos percentuais das regiões Nordeste, Sudeste e Sul.
Considerando -se o desempenho dos alunos, nota -se que os percentuais dos que consideraram a extensão da prova adequada, a partir do segundo quarto de desempenho diferem pouco, em torno de 63 \%.
Para o de pior desempenho ( primeiro quarto ), 60,0 \% consideraram a extensão da prova adequada.
No Gráfico 4.6, pode -se constatar que para três dos quatro quartos de desempenho a proporção de estudantes que consideraram a prova longa ou muito longa, em relação ao tempo total destinado à sua resolução, manteve -se em torno de 31,5 \%, sem diferenças estatisticamente significativas. 4.3 COMPREENSÃO DOS ENUNCIADOS DAS QUESTÕES 4.3.1 Componente de Formação Geral Com relação aos enunciados das questões do Componente de Formação Geral ( Questão 4 ), as opiniões foram positivas, já que 73,6 \% dos alunos avaliados consideraram os enunciados de todas ou da maioria das questões claros e objetivos ( Gráfico 4.7, Gráfico 4.8, e, no Anexo II, a Tabela II.4 ).
Na análise regional, a percentagem de estudantes que avaliaram que todos ou a maioria dos enunciados das questões do Componente de Formação Geral estavam claros e objetivos variou de 72,5 \% na região Sudeste a 76,2 \% na região Nordeste, As diferenças entre as Grandes Regiões não são estatisticamente significativas.
A análise das percepções dos estudantes sobre a clareza e objetividade dos enunciados permite afirmar que todos ou a maioria dos enunciados de questões relativas ao Componente de Formação Geral foram considerados claros e objetivos para a maior parte dos respondentes ( maior do que 72 \% em todas as regiões e maior do que 64 \% para todos os quartos de desempenho ).
Segundo o desempenho, observa -se que a proporção dos que emitiram esta opinião cresce conforme o desempenho aumenta, e que todas as diferenças são estatisticamente significativas.
No quarto superior, a clareza e objetividade de todos ou da maioria dos enunciados das questões foi percebida por 81,4 \% dos alunos e no quarto de desempenho inferior tal avaliação foi emitida por 64,1 \% deles. 4.3.2 Componente de Conhecimento Específico Com relação aos enunciados das questões do Componente de Conhecimento Específico da prova, para 67,4 \% dos estudantes avaliados da Área de Computação a clareza e a objetividade ( Questão 5 ) estavam presentes em todas ou na maioria das questões ( Gráfico 4.9, Gráfico 4.10, e no Anexo II, a Tabela II.5 ).
A maioria dos estudantes de todas as Grandes Regiões brasileiras considerou claros e objetivos todos ou a maioria dos enunciados das questões do Componente de Conhecimento Específico da prova, percentual sempre maior ou igual a 66,0 \%.
Além disso, a diferença entre 71,0 \% ( Nordeste ) e 66,0 \% ( Sudeste ) é estatisticamente significativa.
A proporção de estudantes que consideraram os enunciados das questões claros e objetivos apresenta uma tendência crescente em relação ao aumento de desempenho : mais elevada no quarto superior ( 73,0 \% ) se comparada ao quarto inferior de desempenho ( 62,1 \% ).
E esta diferença é estatisticamente significativa. 4.4 SUFICIÊNCIA DAS INFORMAÇÕES/INSTRUÇÕES FORNECIDAS Ao avaliarem as informações/instruções fornecidas para a resolução das questões ( Questão 6 ), 80,5 \% dos respondentes da Área de Computação de todo o Brasil afirmaram que estas eram até excessivas ou suficientes em todas ou na maioria das questões ( Gráfico 4.11, Gráfico 4.12, e, no Anexo II, a Tabela II.6 ).
Quanto à distribuição de respondentes pelas Grandes Regiões observa -se que a proporção de estudantes que consideraram as informações/instruções fornecidas até excessivas ou suficientes em todas ou na maioria das questões foi sempre superior a 75 \%, chegando a 83,2 \% na região Sul.
Há diferenças estatisticamente significativas entre as duas regiões com percentagens menores ( Sudeste e Norte ) em relação às duas regiões com percentagens mais elevadas ( Sul e Nordeste ) dos que optaram por uma das três primeiras respostas da Questão 6.
Levando -se em conta o desempenho dos participantes, notam -se diferenças estatisticamente significativas entre as opiniões de estudantes de cada quarto de desempenho para o imediatamente superior, excetuando -se do segundo para o terceiro quarto, como mostra o Gráfico 4.12.
O percentual de participantes que avaliaram as informações/instruções como até excessivas ou suficientes em todas ou na maioria das questões foi mais elevado no quarto superior ( 86,0 \% ), percentual superior à média nacional ( 80,5 \% ).
Já no quarto inferior, a suficiência das informações/instruções declarada como até excessiva em todas ou na maioria das questões foi percebida por 75,9 \% dos respondentes. 4.5 DIFICULDADE ENCONTRADA AO RESPONDER À PROVA Perguntados sobre as dificuldades com as quais se depararam ao responder à prova ( Questão 7 ), 21,6 \% dos estudantes apontaram o desconhecimento do conteúdo.
Para 46,5 \%, a forma diferente de abordagem do conteúdo foi indicada como dificuldade.
Já a falta de motivação para fazer a prova foi a dificuldade apontada por 20,6 \% dos respondentes.
Considerando -se todo o Brasil, 7,7 \% dos respondentes afirmaram que não tiveram qualquer tipo de dificuldade para responder à prova ( Tabela II.7 no Anexo II ).
Os Gráficos 4.13 e 4.14 apresentam os percentuais de estudantes que apontaram o desconhecimento do conteúdo como dificuldade percebida ao responder à prova.
Na análise por Grandes Regiões, o percentual de inscritos e presentes que apontaram o desconhecimento do conteúdo como dificuldade ao responder à prova variou de 20,0 \% ( região Nordeste ) a 24,9 \% ( região Centro-Oeste ), diferença estatisticamente significativa.
A forma diferente de abordagem do conteúdo foi a escolha modal dos estudantes, com percentuais que variaram de 42,0 \% ( região Nordeste ) a 50,4 \% ( região Norte ).
O percentual de alunos que citaram a falta de motivação como dificuldade variou de 18,8 \% ( região Norte ) a 25,0 \% ( região Nordeste ).
Os que declararam não ter qualquer dificuldade para responder à prova variaram de 5,5 \% na região Norte a 8,6 \% na Nordeste.
Com relação aos quartos de desempenho, o desconhecimento do conteúdo, foi a opção escolhida por 21,9 \% dos estudantes do quarto superior e 21,0 \% do quarto inferior, evidenciando pouca variação da proporção do que emitiram esta opinião em relação ao desempenho.
A alternativa modal para os alunos, quando agregados pelos quartos de desempenho, foi que a dificuldade encontrada foi causada pela forma diferente de abordagem do conteúdo : 43,3 \% no quarto inferior e 45,3 \% do quarto superior assim o responderam.
Nos quartos intermediários esta proporção foi maior : 47,6 \% no segundo quarto e 49,6 \% no terceiro. 4.6 CONTEÚDOS DAS QUESTÕES OBJETIVAS DA PROVA Ao analisarem os conteúdos das questões objetivas da prova ( Questão 8 ), um percentual muito pequeno dos estudantes avaliados, apenas 8,6 \%, afirmou que não estudou ainda a maioria desses conteúdos ( Gráficos 4.15, Gráfico 4.16, e a Tabelas II.8 no Anexo II ).
A maioria ( 52,1 \% ) afirmou ter estudado e aprendido muitos ou todos os conteúdos.
Na análise por Grande Região, a proporção de respondentes que escolheram a opção não estudou ainda a maioria desses conteúdos, foi pequena.
Observa -se que nas regiões Norte ( 13,4 \% ), Sudeste ( 8,8 \% ) e Centro-Oeste ( 10,7 \% ), apesar de pequenas, as proporções foram maiores do que a média nacional ( 8,6 \% ).
Observa -se diferença estatisticamente significativa entre a região Sudeste e as regiões Sul e Nordeste.
Essas três regiões apresentaram proporções menores dos que optaram por não estudado ainda a maioria desses conteúdos do que as regiões Centro-Oeste e Norte.
Em todas as regiões, a maioria dos presentes afirmou ter estudado e aprendido muitos ou todos os conteúdos, com proporções variando entre 44,2 \% na região Centro-Oeste e 54,8 \% na Sul.
Considerando -se separadamente as opiniões de estudantes dos quatro quartos de desempenho, observa -se que, no quarto inferior, 12,7 \% ofereceram como resposta que não estudou ainda a maioria desses conteúdos, sendo menos de 5 \% os do quarto superior com a mesma resposta.
A diferença entre os alunos que optaram por este motivo de dificuldade nos quartos extremos é estatisticamente significativa.
Tendo em conta o quarto superior, 65,2 \% dos alunos afirmaram ter estudado e aprendido muitos ou todos os conteúdos. 4.7 TEMPO GASTO PARA CONCLUIR A PROVA Ao responderem sobre o tempo de conclusão da prova ( Questão 9 ), mais da metade dos estudantes ( 59,0 \% ) afirmou ter gasto entre duas e quatro horas ( Gráfico 4.17, Gráfico 4.18 e, no Anexo II, a Tabela II.9 ).
Considerando -se as cinco Grandes Regiões brasileiras, os que utilizaram entre duas e quatro horas para finalizar a prova apenas na região Sudeste ( 56,8 \% ) não superaram o percentual nacional ( 59,0 \% ), com diferença estatisticamente significativa em relação às demais regiões.
O percentual dos que dispensaram entre duas e quatro horas para concluir a prova nas outras regiões ficou em torno de 61 \%, como mostra o Gráfico 4.17.
Uma vez sendo analisadas as alternativas escolhidas pelos estudantes que se situam em diferentes quartos de desempenho, observa -se que no quarto inferior uma proporção bem menor de participantes declarou ter gasto entre duas e quatro horas para concluir a prova quando comparados com os do quarto superior, respectivamente 45,2 \% e 70,4 \%.
Observa -se uma tendência crescente com diferenças estatisticamente significativas entre os quartos de desempenho.

% matched lemmas: decrescer
\end{textsample}
